
\paperchapter{\Large}{Compliant Residual DAgger (CR-DAgger)}

Two components:
\begin{itemize}
    \item Compliant Intervention Interface
    \item Compliant Residual Policy
\end{itemize}

\section{Compliant Intervention Interface}

\begin{itemize}
    \item Mainly, they use a standard admittance controller to realise compliant interventions from the expert.
        They designed a hardware addon to allow the human expert to apply forces directly to the robot, and their main takeaway is that we don't need to stop the task while giving demonstrations --- the intervention is done while the robot is still executing its learned policy.

    \item Two benefits of this approach as they claim are that the transitions are smoother, and as a result there is minimal distribution shift between the learned policy and the expert interventions.

    \item During the intervention phase, they sample with a high-frequency, which enables more reactive learning, especially in contact-rich tasks.
\end{itemize}

\section{Compliant Residual Policy}

Their base policy is a diffusion policy with visual and proprioceptive inputs.
They don't have any force inputs in the base, that's included in their Compliant policy, so as to demonstrate force additions can be learned as residuals.

\section{TODO}

\begin{itemize}
    \item Read the HG-DAgger paper (Human-Gated DAgger), which deals with how to add in human intervention in the DAgger framework.

    \item Implement CR-DAgger, as they have mentioned.
\end{itemize}
